\PassOptionsToPackage{table,xcdraw}{xcolor}
\documentclass[10pt, a4paper, landscape]{article}

\usepackage[T1]{fontenc}
\usepackage[utf8]{inputenc}
\usepackage{tgpagella}
\usepackage[spanish, es-noshorthands]{babel}
\usepackage{amsmath, amssymb}
\usepackage[most]{tcolorbox}
\usepackage{tabularx}
\usepackage[margin=1.5cm]{geometry}
\usepackage{fancyhdr}

\pagestyle{fancy}
\fancyhf{}
\setlength{\headheight}{16pt}
\lhead{\textbf{UCB Sede Tarija} | Ing. Mecatrónica}
\rhead{IMT-221 | Matriz de Progreso de Laboratorio}
\renewcommand{\headrulewidth}{0.4pt}
\definecolor{ucbblue}{RGB}{0, 51, 102}

% Definición de columnas robustas: evitan el colapso de tabularx (Emergency Stop) 
% forzando el ajuste de texto y el salto de línea.
\newcolumntype{C}[1]{>{\centering\arraybackslash\hspace{0pt}}p{#1}}
\newcolumntype{L}{>{\raggedright\arraybackslash\hspace{0pt}}X}

\begin{document}

\begin{tcolorbox}[colback=ucbblue!5, colframe=ucbblue, title=\textbf{Matriz de Progreso Dual-Track (Hardware y Simulación) — Hito 2}]
    Esta matriz monitorea el avance equitativo de todos los grupos. \textbf{Track A} llena campos físicos; \textbf{Track B} llena campos analíticos y simulados, dejando mediciones físicas explícitamente pendientes[cite: 18].
\end{tcolorbox}

\vspace{0.3cm}

{\footnotesize % Reducimos la fuente para acomodar 13 columnas holgadamente
\renewcommand{\arraystretch}{2}
\begin{tabularx}{\textwidth}{|C{0.6cm}|C{1.1cm}|L|C{1.3cm}|C{1.5cm}|C{1.3cm}|C{1.5cm}|C{1.5cm}|C{1.5cm}|C{1.5cm}|C{1.3cm}|L|L|}
\hline
\rowcolor{ucbblue!20}
\textbf{Grp} & \textbf{BOM OK?} & \textbf{Faltantes} & \textbf{Pinout OK?} & \textbf{$I_T$ analít.} & \textbf{$I_T$ sim.} & \textbf{$I_T$ medido} & \textbf{Q sim.} & \textbf{Q medido} & \textbf{$A_d$ pred.} & \textbf{Estado} & \textbf{Problema} & \textbf{Próxima Acción} \\ \hline
\textbf{1} & & & & & & & & & & & & \\ \hline
\textbf{2} & & & & & & & & & & & & \\ \hline
\textbf{3} & & & & & & & & & & & & \\ \hline
\textbf{4} & & & & & & & & & & & & \\ \hline
\textbf{5} & & & & & & & & & & & & \\ \hline
\textbf{6} & & & & & & & & & & & & \\ \hline
\textbf{7} & & & & & & & & & & & & \\ \hline
\textbf{8} & & & & & & & & & & & & \\ \hline
\end{tabularx}
}

\vspace{0.8cm}
\textbf{Criterios de Evaluación y Exigencia Matemática por Columna[cite: 18]:}
\begin{itemize}
    \item \textbf{BOM/Pinout:} Validación documental en el datasheet. No se acepta "asumir" el pinout del 2N3904.
    \item \textbf{$I_T$ analítico:} Exigir la derivación matemática KVL para el espejo de corriente: 
    \begin{equation*}
        I_{REF} = \frac{0\,\text{V} - (V_{EE} + 0.7\,\text{V})}{R_{REF}}
    \end{equation*}
    Y la corrección teórica considerando un $\beta$ finito para obtener la corriente de cola: $I_T \approx \frac{I_{REF}}{1+2/\beta}$[cite: 18].
    \item \textbf{Puntos Q (simulado y medido):} Los datos deben mostrar coherencia con KVL/KCL y operación activa. Exigir explícitamente la verificación de la inecuación para ambas ramas: 
    \begin{equation*}
        V_{BC} = V_B - V_C < 0\,\text{V} \quad \text{[cite: 18]}
    \end{equation*}
    \item \textbf{$A_d$ predicho:} Debe ser calculado matemáticamente a partir del Punto Q real (medido o simulado por el estudiante), \textbf{NO} a partir del valor nominal ideal. Ecuación exigida para validación: 
    \begin{equation*}
        A_{d,\text{pred}} = g_{m,\text{real}} R_{C,\text{real}} = \left(\frac{I_{CQ,\text{real}}}{V_T}\right) R_{C,\text{real}} \quad \text{[cite: 18]}
    \end{equation*}
\end{itemize}

\end{document}
